\subsection{局部平凡纤维空间}

设 B 为拓扑空间。若 F 是拓扑空间，我们称 B-空间 \((B \times F, pr_{1})\) 为以 F 为纤维型的平凡 B-纤维空间。

设 E 为 B-空间。若存在拓扑空间 F 以及 B-同构 \(u: E \rightarrow B \times F\)，则称 B-空间 E 为可平凡化的 B-纤维空间，并称 u 为它的一个以 F 为纤维型的平凡化。

定义 1. --- 设 B 为拓扑空间。设 E 为 B-空间，p 为其投影。若 B 的每一点都有邻域 V 使 \((\overline{p}^{1}(V), p_{V})\) 是可平凡化的 V-纤维空间，则称 E 为局部平凡 B-纤维空间。

除了说“局部平凡 B-纤维空间”之外，也可说“以 B 为基的局部平凡纤维空间”。若 E 是局部平凡 B-纤维空间，有时也说 \((E,B,p)\) 是一个局部平凡纤维化，或者，滥用语言，说 p 是一个局部平凡纤维化。也称 B-空间 \((E,p)\) 在 B 的子集 A 上是可平凡化的，如果 A-空间 \(E_{A}\) （由 \((E,p)\) 在 A 上诱导者）是可平凡化的纤维空间。

设 E 为局部平凡 B-纤维空间，用 p 表示其投影。使纤维 \(\overline{p}^{1}(a)\) 为空（相应地，非空）的 B 中的点 a 所成的集是开的。于是 p 的像是 B 中既开又闭的子集。

设 F 为拓扑空间。若 E 的所有纤维都同胚于 F，则称 E 为以 F 为纤维型的局部平凡 B-纤维空间。

注. --- 1) 设 (E, p) 是 B 上的局部平凡纤维丛，并设 A 是 B 的子集。由 E 通过过渡到子空间而得的 A-空间 \(\mathrm{E}_{\mathrm{A}} = (\stackrel{-1}{p}(\mathrm{A}), p_{\mathrm{A}})\) 是局部平凡纤维空间。若 E 是可平凡化的，则 \(E_{A}\) 亦然。

事实上，对 A 的每一点 \(a\)，存在 \(a\) 在 B 中的一个开邻域 U、一个拓扑空间 F 以及一个 U-同构 \(g\colon \overline{p}^{1}(\mathrm{U})\to \mathrm{U}\times \mathrm{F}\) 。映射 \(g\) 诱导出一个 \((\mathrm{A}\cap \mathrm{U})\) -同构，它把 \(\overline{p}_{\mathrm{A}}^{-1}(\mathrm{A}\cap \mathrm{U})\)

映到 \((\mathrm{A}\cap\mathrm{U})\times\mathrm{F}\) 上，这就证明了 \(E_{A}\) 是局部平凡 A-纤维空间。

\begin{enumerate}
\def\labelenumi{\arabic{enumi})}
\setcounter{enumi}{1}
\tightlist
\item
  设 \((\mathrm{E}, p)\) 为 B-空间，并设 \((\mathrm{E}_{i})_{i \in \mathrm{I}}\) 为由开集组成的 E 的一个分划。
\end{enumerate}

若 B-空间 \((\mathrm{E}_{i}, p \mid \mathrm{E}_{i})\) 中的每一个都是可平凡化的纤维空间，则 E 是可平凡化的 B-纤维空间。事实上，对每个 \(i \in I\)，设 \(F_{i}\) 为使 B-空间 \((\mathrm{E}_{i}, p \mid \mathrm{E}_{i})\) 与 \((\mathrm{B} \times \mathrm{F}_{i}, \mathrm{pr}_{1})\) 同构的拓扑空间。于是 B-空间 \((\mathrm{E}, p)\) 同构于 \((\mathrm{B} \times \mathrm{F}, \mathrm{pr}_{1})\)，其中 F 是族 \((\mathrm{F}_{i})_{i \in I}\) 的拓扑和空间。

设集 I 是有限的。若 B-空间 \((\mathrm{E}_{i}, p \mid \mathrm{E}_{i})\) 中的每一个都是局部平凡 B-纤维空间，则 B-空间 E 亦然。事实上，由于集 I 是有限的，B 的每一点都有邻域 U，使得诸 B-纤维空间 \(E_{i}\) 在 U 上都是可平凡化的。由前所证，U-空间 \((\bar{p}^{1}(\mathrm{U}), p_{\mathrm{U}})\) 于是是可平凡化的纤维丛。

